% --- 自動計數器設定 ---
\newtotcounter{qcount} \regtotcounter{qcount}
\newtotcounter{sccount} \regtotcounter{sccount} 
\newtotcounter{mccount} \regtotcounter{mccount} 

% --- 基礎縮排設定 ---
\newcommand{\choiceindent}{0.85cm}

% --- 1. 防分頁容器 (qblock) ---
\newenvironment{qblock}{%
    \par\vspace{0.5em}
    \noindent\begin{minipage}{\textwidth} % 確保整題不分頁
}{%
    \end{minipage}\par\vspace{1em}
}

% --- 2. 題幹指令 (純計數) ---
\newcommand{\schoices}[7]{% 
    \begin{qblock}
        \stepcounter{qcount}\stepcounter{sccount}%
        \begin{enumerate}[label=\textbf{\theqcount.}, leftmargin=\choiceindent, labelsep=0.2cm, nosep]
            \item #1
        \end{enumerate}\vspace{0.2em}
        \ifstrequal{#2}{h}{
            \noindent\hspace*{\choiceindent}% 
            (A) #3 \hspace{0.5cm} (B) #4 \hspace{0.5cm} (C) #5 \hspace{0.5cm} (D) #6 \hspace{0.5cm} (E) #7 \par
        }
        {
            \begin{enumerate}[(A), topsep=0.3em, itemsep=0.5em, parsep=0pt, leftmargin=*, labelindent=\choiceindent, labelsep=0.25cm, align=left] 
                \item #3 \item #4 \item #5 \item #6 \item #7
            \end{enumerate}
        }
    \end{qblock}
}

\newcommand{\mchoices}[8]{% 
    \begin{qblock}
        \stepcounter{qcount}\stepcounter{mccount}%
        \begin{enumerate}[label=\textbf{\theqcount.}, leftmargin=\choiceindent, labelsep=0.2cm, nosep]
            \item #1 \textbf{(應選 #2 項)}
        \end{enumerate}\vspace{0.2em}
        \ifstrequal{#3}{h}{
            \noindent\hspace*{\choiceindent}% 
            (A) #4 \hspace{0.5cm} (B) #5 \hspace{0.5cm} (C) #6 \hspace{0.5cm} (D) #7 \hspace{0.5cm} (E) #8 \par
        }
        {
            \begin{enumerate}[(A), topsep=0.3em, itemsep=0.5em, parsep=0pt, leftmargin=*, labelindent=\choiceindent, labelsep=0.25cm, align=left] 
                \item #4 \item #5 \item #6 \item #7 \item #8
            \end{enumerate}
        }
    \end{qblock}
}

% --- 側邊圖片包裝指令 ---
% #1: 題目內容 (填入 \schoices{...} 或 \mchoices{...})
% #2: 圖片檔名
% #3: 圖片佔頁面寬度的比例 (例如 0.1)
\newcommand{\sideimg}[3]{%
    \begin{qblock} % 使用你之前的防分頁環境
        \noindent
        % 左側：自動計算剩餘寬度 (總預算 0.91 扣掉圖片寬度)
        \begin{minipage}[t]{\dimexpr 0.91\textwidth - #3\textwidth \relax}
            % 這裡必須要把 \schoices 原本內建的 before=\vspace 抵消掉，
            % 否則圖片跟第一行字會對不齊，所以稍微修正間距
            \vspace{-0.5em} 
            #1
        \end{minipage}
        \hfill 
        % 右側：放圖片
        \begin{minipage}[t]{#3\textwidth}
            \includegraphics[width=\textwidth, valign=t]{#2}
        \end{minipage}
    \end{qblock}
}

\newcommand{\singlechoice}[1]{% 
    \stepcounter{qcount}\stepcounter{sccount}%
    \begin{enumerate}[label=\textbf{\theqcount.}, leftmargin=\choiceindent, labelsep=0.2cm, nosep]
        \item #1
    \end{enumerate}\vspace{0.2em}
}

\newcommand{\multiplechoice}[2]{% 
    \stepcounter{qcount}\stepcounter{mccount}%
    \begin{enumerate}[label=\textbf{\theqcount.}, leftmargin=\choiceindent, labelsep=0.2cm, nosep]
        \item #1 \textbf{(應選 #2 項)}
    \end{enumerate}\vspace{0.2em}
}

% --- 3. 自動比例圖片 (sidefig) ---
\newcommand{\sidefig}[3]{% {1:題幹文字}{2:圖片檔名}{3:圖片比例}
    \begin{minipage}[t]{\dimexpr 0.91\textwidth - #3\textwidth \relax}
        #1
    \end{minipage}\hfill 
    \begin{minipage}[t]{#3\textwidth}
        \includegraphics[width=\textwidth, valign=t]{#2}%
    \end{minipage}
}

% --- 4. 選項排版 ---
\newcommand{\hchoices}[5]{% 水平排列
    \noindent\hspace*{\choiceindent}% 
    (A) #1 \hfill (B) #2 \hfill (C) #3 \hfill (D) #4 \hfill (E) #5 \par
}

\newcommand{\vchoices}[5]{% 垂直排列
    \begin{enumerate}[(A), nosep, leftmargin=*, labelindent=\choiceindent, labelsep=0.25cm, align=left] 
        \item #1 \item #2 \item #3 \item #4 \item #5
    \end{enumerate}
}